\section{One machine for all exponents}\label{sec:single-machine}

Theorem~\ref{thm:integer} gives a multiplier for each exponent above
$1/2$. Choosing a better exponent may change its program, the number
of tapes and the constant in its space bound. To prove
Theorem~\ref{thm:uniform-main} we construct one machine with a fixed
number of heads and a fixed linear bound on the cells they visit.
It runs a sequence of correct multiplication programs under time
limits. Alongside each program it tries several operand lengths,
so that even a program with a large space constant can succeed
within the common space bound.

\subsection{Constructing the candidate programs}

For $j\ge1$ take the target exponent
$\theta_j=1/2+1/(j+2)$ in Theorem~\ref{thm:integer}.
Choose rational intermediate exponents with strict slack.
We need a generator that, given $j$, constructs the corresponding
finite templates and determines how much precision to use at any
input length.

\paragraph{Finding the templates.}
The repair, Fourier and binary templates of
\cite[Sections~\ext{F}{sec:repair} and~\ext{F}{sec:prime-kernel},
including Section~\ext{F}{sec:binary}]{oaF}, and the interchange
template of Section~\ref{sec:interchange}, are built from finite
sets, integer polynomials and linear algebra over explicitly
presented rings or fields. Failing offsets can be enumerated,
rational bases computed by elimination, and radix primes tested
against the finite lists of numerators and denominators. Thus each
completed construction supplies a definite list of instructions,
widths and recursive calls. Enumerate the integer template
parameters, interleaving the searches if necessary. Accept a
construction when its call counts meet the strict inequalities in
the two recursions. Raising to integer powers decides a comparison
between a positive rational and a rational power of a positive
integer. The existence proofs guarantee termination for every $j$.

For a given input length the generated program tries the trial
divisions and checks the volume, coprimality and padding conditions
of Equation~\eqref{eq:integer-box}. It verifies the space for
borrowed coordinates and for row padding in a Walsh layer, and the
digit-width bound $b\le4k$ used in the precision estimates. Here
$N=2^k$ is the padded number of numerical records. The direct splits
of Section~\ref{sec:permutations} apply at small address widths.
Use schoolbook multiplication at degenerate lengths or if any
required shape test fails. For a fixed $j$ the asymptotic estimates
prove that the fast construction applies at every sufficiently large
length. The program can therefore handle short inputs without
having to compute a cutoff from a prime-gap estimate.

\paragraph{Certifying enough precision.}
We describe how the program obtains its working widths from the
finite description and the accepted data shape. At a pointwise
instruction in the Walsh template, the sum of the absolute matrix
entries bounds the operator norm on its live row registers. Include
the identities on retained registers and take the maximum over the
finite coefficient table. Copies and merges have bounds from their
finite multiplicities. A permutation or a completed normalized child
transform has norm one. Choose an integer $C_j\ge2$ that bounds the
product of all the one-level instruction bounds, each enlarged to
at least one, including the base case and the direct butterfly.

For template arity $M\ge2$ the recursion depth is at most
$d=1+\lceil\log_M(1+k)\rceil$. To propagate an error between two
internal instructions, follow their paths to their common ancestor.
Each path costs at most $C_j^d$ in gain, and the instructions at the
ancestor cost at most $C_j$. The same $C_j^{2d+1}$ bound works
across consecutive pieces since each completed piece is unitary.
The number of array instructions is polynomial in $1+k$ with
computable exponent, by fixed branching and logarithmic depth.
Enumerating these instructions takes $\operatorname{poly}_j(1+k)$
setup time. It gives explicit polynomial bounds on both individual
gains and the sum of gains for propagated faults.

At the boundaries of the transforms use a norm cap at least two
above the ideal Euclidean norms obtained from the digit bound,
unitarity, the final scale $\sqrt S$ and
$\|v\odot w\|_2\le\|v\|_2\|w\|_2$. The cap is $2^{O(k+1)}$ on the
accepted shapes. For a fractional grid $2^{-P}$, an error in a
coefficient costs its absolute size times the operand norm.
Coordinate rounding costs at most $2\sqrt{N'}\,2^{-P}$ in a stream
of $N'=O_j(N)$ padded records. Within a layer use the boundary cap
times the computed gain, allowing one extra unit of internal
Euclidean error. The roots and square roots in Section~\ref{sec:integers}
give certified local tables. In a product of at most $k$ phases,
each with coefficient and rounding errors at most $2^{-P}$, the
total coefficient error is bounded by
$2k(1+2^{-P})^k2^{-P}$. Taking $2^P\ge2k$ makes the parenthesized
power at most two.

Multiply each local error bound by the corresponding propagation
gain and sum. At the boundary of a completed linear layer the
incoming discrepancy carries through with unit gain. At the one
coordinatewise product, twice the norm cap bounds the gain from
the larger of the two input discrepancies. Include the final
scaling by $\sqrt S$. These operations give an integer majorant
$H_j(k)=2^{O_j(k+1)}$ for the coefficient of $2^{-P}$ in every
error bound, including the internal errors of a layer executed on
its represented input. Take
\[
                     2^P>8H_j(k)+2k
\]
and allow enough fractional bits to encode the fixed dyadic
coefficients exactly. Take the integer width from the computed
magnitude bounds, including the extra unit. Induction through
the instructions now preserves the caps. Each internal discrepancy
from the exact layer on its represented input is less than one,
and the discrepancies at completed boundaries are less than $1/4$.
Rounding the final coefficients gives the exact product. The
working widths are $O_j(1+k)$, with the polynomial setup and
time orders of Section~\ref{sec:integers}. The calculation uses
actual instruction counts, so it also certifies short accepted
shapes.

\subsection{Using a common number of tapes}

At a template node the live streams are aligned arrays. Encode the
finitely many roles of a fixed candidate as tracks of one sequence
of records. To replace a track, scan its data to scratch, compute
its new entries and merge them back by a scan. Several tracks can
be accumulated successively. To split into finitely many residue
classes, scan out one class at a time and mark its positions for
the return merge. This serializes the fixed fanout in the radix
streams. Elementary scalings and earlier-field controlled
translations in an ordered affine map use the bounded number of
fields, counters and rotation streams of a single elementary
operation. Any fixed split or merge inside it can be serialized
in the same way. The comparisons for the spectral table in
Section~\ref{subsec:swap-spectrum} produce a binary mask on the
role tape, one comparison at a time.

The resulting scans at a node cost $O_j(V)$ for its local bit
volume $V$. For a columnwise Walsh address selector, extract the
finitely many address slots in separate passes. Merge each slot
with the successive column tuples to append its bit to each
tuple. A sequential pass then computes the selector. This
takes $O_j(k)$ per address, within the record-processing budget.

Before a recursive call, extract its single input stream. Park
the other live tracks and the return descriptor at the end of a
stack tape. A full-volume reserve descent of interchange parks
only the descriptor. Restore the frame from that same end on
return and merge the child's output. All the extra scans belong
to the parent and occur only a fixed number of times there. In
particular the child receives its original volume, so the
coefficients of the recursive costs in the interchange and Walsh
bounds are preserved. A symbol of a fixed candidate alphabet can
be carried as a fixed-length binary block during interchange.

Sorting, table iteration, arithmetic in a record, reversal and
output use a bounded number of streams. The candidates share the
fixed short-integer multiplier. Together with the track processors
and end-access stacks these routines use a common number $K$ of
physical tapes over one fixed alphabet. The preceding construction
therefore generates correct $K$-tape programs $A_j$ with
\[
          T_j(u)=O_j\bigl(u(1+\log u)^{\theta_j}\bigr)
\]
time on two $u$-bit operands. A fixed interpreter can simulate
each candidate with one coded work tape per head and one tape for
the transition list. For fixed $j$, scanning that list and updating
the tapes take constant time per step. Generating the program is a
constant setup cost for that $j$.

Let $c_j(u+1)$ bound the physical cells visited by a complete
interpreted run, including generation. Such a constant $c_j$ exists
by the fixed-template space bounds. The serialization uses a
constant number of copies of each stream for fixed $j$, stores the
same parent frames, and adds a finite program description. The
schedule below never needs the value of $c_j$.

\subsection{Scheduling within linear space}

Mark a distance $n+O(1)$ on either side of the origin of each
candidate work tape, once in $O(n)$ time. Abort a run if it reaches
beyond its marked region, even during program generation. Maintain
the original inputs, the tentative answer and the scheduling data
on separate tapes of fixed linear capacity.

An attempt chooses a program index $j\ge1$ and a chunk count
$r=2^k$, $k\ge0$, with $r\le n$. Skip larger counts. With
$u=\lceil n/r\rceil$, left-pad each input to $r$ digits in
base $2^u$. Generate $A_j$ and use it for the $r^2$ digit
products, one at a time. Accumulate the shifted products in
$2ru\le4n$ bits and discard the leading padding on output.
Only an attempt that completes all these operations returns a
result. Clear and reuse the visited candidate work area between
calls. For fixed $r$, gathering the operands by sequential
scans, positioning the additions and propagating carries cost
$O_r(n)$ in total. The loop counts and shifts use $O(1+\log n)$
bits. Unary position marks or amortized binary counters suffice
during the scans.

For each $j$, some fixed power of two $r>2c_j$ makes a digit
multiplication fit within the marked tapes at every sufficiently
large $n$, including its setup. That attempt runs to completion
in time
\[
 O_{j,r}\bigl(n+r^2u(1+\log(2u))^{\theta_j}\bigr)
              =O_{j,r}\bigl(n(1+\log n)^{\theta_j}\bigr).
\]
All attempts, including ones that time out or cross a mark, share
the same absolute $O(n)$ bound on visited cells.

Use a weighted schedule as in universal search \cite{Levin}.
Number the pairs $(j,k)$ in computable diagonal order, reserving
index $1$ for a schoolbook multiplier with its own fixed linear
space allowance. In round $s=1,2,\ldots$, restart the first $s$
attempts in turn. Give index $i$ at most $2^{s-i}$ actual steps
and stop at the first completed output. Decoding the pair, generating
its candidate, copying operands and accumulating the tentative
result all count against the budget. A clock for $2^a$ steps is a
binary counter with its head at the low end. Marking the highest
nonzero bit detects exhaustion. Over $v$ ticks the decrements and
head returns cost $O(a+v+1)$ in total by summing the borrow
lengths, including the initial decrement. At an abort or timeout,
erase the visited scratch and tentative output and rewind the
heads. Frontier marks track the visited intervals, so cleanup
costs time proportional to the work of the attempt and its clock
setup. The outer space marks stay in place across rounds.

The whole of round $s$ costs $O(2^s)$, including the polynomial
overhead in $s$ for its clocks and control. Schoolbook multiplication
terminates in $O(n^2)$ steps, so the last round has $s=O(1+\log n)$.
All clocks and scheduling data fit within the claimed space bound.
For a fixed successful attempt of index $i$ taking $t$ steps,
round $i+\lceil\log_2(t+1)\rceil$ allows it to finish. The sum of
all round costs up to that point is $O(2^i(t+1))$, in addition
to the one initial $O(n)$ scan. For every $j$ a fixed index has
time $O_j(n(1+\log n)^{\theta_j})$ at sufficiently large $n$.
Since $\theta_j\downarrow1/2$, this one machine has the time bound
$n(1+\log n)^{1/2+o(1)}$ and visits $O(n)$ cells, proving
Theorem~\ref{thm:uniform-main}.

\subsection{Quotients and square roots}
\label{sec:exact-arithmetic-applications}

\begin{corollary}[Exact division and integer square root]\label{cor:near-root-applications}
On the finite-alphabet multitape model of Theorem~\ref{thm:uniform-main},
each of the following tasks has a single deterministic machine running
in worst-case time $m(1+\log m)^{1/2+o(1)}$ with $O(m)$ visited cells.
For two $m$-bit integers $a,b$ with $b>0$, division returns
$\lfloor a/b\rfloor$ and $a\bmod b$, each in $m$ bits. For an
$m$-bit integer $a$, square root returns $\lfloor\sqrt a\rfloor$
in $m$ bits. Leading zeros are allowed.
\end{corollary}

\begin{proof}
Use Newton's classical reductions \cite{Brent1976}, as for the
earlier tape multiplier \cite[Section~10 and Corollary~2]{Tape}.
At each stage the number of accurate bits doubles. The same
machine from Theorem~\ref{thm:uniform-main} performs all the
multiplications. For each fixed $\eta>0$, its cost on $u$-bit
inputs is $O_\eta(B_\eta(u))$, where
$B_\eta(u)=u(1+\log u)^{1/2+\eta}$. Neither Newton routine
takes $\eta$ as input.

\emph{Reciprocal.} Normalize the divisor as $b=2^h c$ with
$1\le c<2$. The update $x\gets x(2-cx)$ squares the residual
$1-cx$ in exact arithmetic. Truncate $c$ to $t$ fractional bits
and round down each of the two products to this precision. When
$|1-cx|\le1/2$, the new absolute residual is at most
\[
                       (1-cx)^2+6\cdot2^{-t}.
\]
The truncation contributes at most $cx^2\,2^{-t}$ and the product
roundings at most $(cx+c)2^{-t}$. Here $cx^2\le9/4$ and
$cx+c\le7/2$. The seed $x=5/8$ has residual at most $3/8$,
so one step at $t=6$ gives residual below $1/4$. Given residual
at most $2^{-j-1}$, take $j'=\min(2j,P)$ and $t=j'+5$ for target
precision $P$. The new residual is at most
$2^{-2j-2}+6\cdot2^{-j'-5}<2^{-j'-1}$.

\emph{Square root.} For a positive radicand write $a=2^{2e}c$,
$1\le c<4$. Start with $x=3/2$ and update $x\gets(x+c/x)/2$.
Its exact error is $(x-\sqrt c)^2/(2x)$. In the range
$1/2\le x\le5/2$ the reciprocal routine, after a binary shift
to normalize $x$, gives $1/x$ to absolute error at most $2^{-t}$
at target $t$. Indeed the residual is at most $2^{-t-1}$ and
$1/x\le2$. Truncating $c$, multiplying by the approximate
reciprocal and rounding the average change the update by at most
$(9/2)2^{-t}$. The quotient error is at most $(1+4+2)2^{-t}$
before averaging. At $t=6$ the first step has error at most
$1/12+(9/2)2^{-6}<1/4$. The same capped doubling schedule now
preserves $|x-\sqrt c|\le2^{-j-1}$ and keeps $x$ in $[3/4,9/4]$.

Use $P=2m+4$ in each routine, handling a zero dividend or
radicand by a scan. Multiplying the normalized reciprocal by $a$
and applying the binary scale gives an approximation to $a/b$
with error less than $1/4$. The rescaled root also has error
less than $1/4$. Take the floor and correct by at most one,
testing the candidate against the input by exact multiplication
and comparison. For division this also supplies the remainder.

Make a copy of the normalized input in one scan so that a stage
of precision $t$ reads only its own prefix. A reciprocal stage
uses two $O(t)$-bit products, and a root stage uses one reciprocal
computation and one additional product of that size. The sum of
the precisions is $O(P)$. Since $B_\eta(t)/t$ is nondecreasing,
summing the reciprocal stages gives $O_\eta(B_\eta(P))$ time.
Applying this bound to the internal reciprocal calls and summing
the root stages gives the same time. The initial scans, final
products, corrections and linear output padding fit this budget.
The fixed scratch tapes are reused at bounded call depth, and the
total precision allocated along a call is $O(m)$. Rewinding and
clearing each scratch area cost at most the time spent visiting
it. The two routines accordingly visit $O(m)$ cells.
\end{proof}

\begin{corollary}[Lower bounds for division and integer square root]\label{cor:division-root-lower}
Assume variable-count Weak ETH. Every fixed deterministic multitape
machine computing either the quotient of two $m$-bit integers with
positive divisor or the floor square root of an $m$-bit integer takes
worst-case time $\Omega(m\sqrt{\log m})$ at every sufficiently large
$m$. The constant and threshold may depend on the machine.
\end{corollary}

\begin{proof}
Brent proved linear equivalences among the corresponding operations
at multiple floating-point precision \cite[Theorem~3.1]{BrentMultiplePrecision}.
Here shifts and additions give reductions to exact integer operations
at a prescribed length. Fix $0\le r<20$. On $k$-bit inputs $x,y$ set
$m=20k+r$ and $B=2^{3k+7}$, padding every argument of a division or
square-root call to exactly $m$ bits.

For division, one call yields
\[
 \left\lfloor\frac{xB^2}{B-y}\right\rfloor=xB+xy,
 \qquad 0\le\frac{xy^2}{B-y}<1 .
\]
Subtracting the shift $xB$ gives the product. For square root, let
$t$ be one of $x,y,x+y$. A call on $B^4+4tB^3$ returns
$W=B^2+2Bt-2t^2$. Indeed $0\le t<2^{k+1}$ gives $B>8t^3$ and
$W\ge B^2$, so
\[
 0\le B^4+4tB^3-W^2=8Bt^3-4t^4<2W+1 .
\]
Subtracting $W$ from $B^2+2Bt$ and shifting recovers $t^2$.
The three squares yield $2xy=(x+y)^2-x^2-y^2$.

For $k\ge4$ the largest square-root argument is less than
$2B^4=2^{12k+29}\le2^m$, and both division arguments also fit.
Constructing these arguments and extracting the result uses
$O(m)$ tape steps. If $D(m)$ denotes the worst-case time of the
given machine, each fixed $r$ thus gives a multiplier with time
$O(D(20k+r)+k)$, using schoolbook multiplication for $k<4$.
Theorem~\ref{thm:lower} applied to these twenty fixed multipliers
gives the assertion on every residue class. Taking the smallest
constant and the largest threshold gives the bound at every
sufficiently large $m$.
\end{proof}
